\documentclass[11pt]{article}
\usepackage{icmlko}
\begin{document}

\title{규약의 폭이 낡았다 --- 미결정 폭이 모든 결정을 팝업 65행에 넘긴다}
\author{월드모델 실험실}
\date{노트 366 · 2026년 8월 1일}
\maketitle

\begin{abstract}
노트 365가 만화 258행 앞에서 막혔다 --- 넣으면 \texttt{media\_push} 채움이
자리 기준이라 축이 조용히 사라진다. 우회로를 물었다: \emph{그 축이 아무것도
안 벌면 그냥 끄고 넣으면 된다.} 쟀다. \textbf{번다} --- 끄면 판
$-0.0025$(SD 0.0026, \textbf{3/12})다. 우회로는 막혔고 채움을 제대로 다시
만드는 수밖에 없다(재료는 긁는 중이다).

\textbf{그런데 그 실험이 더 큰 것을 드러냈다.} 끄면 \emph{팝업이}
$+0.0227$(\textbf{11/12}) 오른다. 노트 275의 규칙대로면 판 $-0.0025$ 는
미결정 폭 $0.010$ 안이고 어느 도메인도 $|t|>2.77$ 로 안 나빠지므로
\textbf{배포가 정하고, 배포는 ``버려라''고 한다} --- 판에서 일하고 있는
축을. 이상해서 세어 봤다. \textbf{이번 창의 짝 실험 여덟 개가 전부 폭
안이고, 제일 큰 것도 폭의 41\%다.} 그러니 그 규칙은 사실상
\textbf{모든 결정을 팝업에 넘긴다} --- 노트 351 · 356 · 358 · 362, 그리고
여기까지 \textbf{다섯 번 연속}이다. 팝업은 유보 65행에 95\% 구간
$[0.223,\,0.661]$(노트 352)인 이 판에서 제일 안 재지는 수다.
\textbf{이번에는 안 버린다} --- 판 부호가 9/12 로 서 있고, 그것은 노트
351이 채택할 때 쓴 것과 같은 증거 세기다. 폭을 고치는 것은 다음 노트로
미리 등록한다. \emph{내 결정을 바꾸는 노트에서 그 결정의 규칙을 같이
고치지 않는다.}
\end{abstract}

\section{우회로는 막혔다}

\begin{center}
\begin{tabular}{lrrr}
\hline
만화 \texttt{media\_push} 를 끄면 & 차 & SD & 부호 \\
\hline
\textbf{판} & $-0.0025$ & 0.0026 & \textbf{3/12} \\
게임 & $-0.0162$ & 0.0094 & 0/12 \\
도서 & $-0.0259$ & 0.0316 & 2/12 \\
\textbf{팝업}(대표) & $\mathbf{+0.0227}$ & 0.0242 & \textbf{11/12} \\
아이돌 & $+0.0187$ & 0.0383 & 7/12 \\
\hline
\end{tabular}
\end{center}

노트 187의 검사를 먼저 통과했다 --- 마스크가 $1.000 \to 0.000$ 으로
실제로 꺼졌고 판이 움직였다. \textbf{만화가 채점도 안 되는데 그 축이 판을
$0.0025$ 만큼 받치고 있다.} 남에게 주는 값이다.

\section{그런데 규칙이 버리라고 한다}

노트 275: ① 판이 미결정 폭($0.010$) 안이고 ② 어느 도메인도 $|t|>2.77$ 로
안 나빠지면 \textbf{배포 쪽 이득이 정한다.}

여기서 ① $|-0.0025| < 0.010$ ② 제일 나쁜 도서가 $-0.0259/0.0316$ 로
$|t|=0.82$. \textbf{둘 다 통과한다.} 그러면 배포가 정하고 팝업이
$+0.0227$(11/12)이니 \emph{버리라}는 답이 된다 --- 판에서 일하는 축을.

\section{세어 보니 폭이 낡았다}

\begin{center}\includegraphics[width=\linewidth]{figs/band.pdf}\end{center}

\begin{center}
\begin{tabular}{lrrrr}
\hline
실험 & 판 차 & SD & 부호 & $|$차$|/0.010$ \\
\hline
만화 없음(361)        & $+0.0041$ & 0.0062 & 8/12 & \textbf{0.41} \\
만화 10\%(362)        & $+0.0032$ & 0.0066 & 9/12 & 0.32 \\
만화 \texttt{media\_push}(366) & $-0.0025$ & 0.0026 & 3/12 & 0.25 \\
만화 원핫 뺌(362)     & $-0.0018$ & 0.0017 & 1/12 & 0.18 \\
팝업 등급 풀기(358)   & $+0.0010$ & 0.0048 & 6/12 & 0.10 \\
만화 절반(361)        & $-0.0008$ & 0.0046 & 6/12 & 0.08 \\
아이돌 위키 뺌(359)   & $-0.0007$ & 0.0016 & 3/12 & 0.07 \\
팝업 전용 축(359)     & $+0.0001$ & 0.0001 & 2/12 & 0.01 \\
\hline
\end{tabular}
\end{center}

\textbf{여덟 개가 전부 폭 안이다.} 제일 큰 것이 41\%이고 절반은 10\%
아래다. 노트 275가 $0.010$ 을 고를 때(노트 245의 미결정 폭)는 판이 그만큼
움직이던 시절이었다. \textbf{지금은 아무것도 그 폭에 안 닿는다} ---
그러니 그 조항은 ``판이 못 정하는 드문 경우''가 아니라 \textbf{기본값}이
됐다.

\section{그래서 팝업이 다 정했다}

\begin{center}
\begin{tabular}{llrrl}
\hline
노트 & 처치 & 판 & 팝업 & 결정 \\
\hline
351 & 검색·위키 컷 & $+0.0002$ (7/12) & $+0.0168$ (11/12) & 넣음 \\
356 & 만화 뺌      & $+0.0058$ (10/12) & $-0.0340$ (2/12) & 안 뺌 \\
358 & 등급 풀기    & $+0.0010$ (6/12) & $+0.0140$ (9/12) & 넣음 \\
362 & 만화 10\%    & $+0.0032$ (9/12) & $-0.0315$ (3/12) & 안 넣음 \\
366 & \texttt{media\_push} & $-0.0025$ (3/12) & $+0.0227$ (11/12) & \textbf{안 버림} \\
\hline
\end{tabular}
\end{center}

\textbf{다섯 번 연속으로 판이 아니라 팝업이 정했다.} 그리고 팝업은 유보
65행 · 95\% 구간 $[0.223,\,0.661]$(노트 352)로 \textbf{이 판에서 제일 안
재지는 수}다. 노트 360이 ``팝업에 필요한 것은 센 행''이라 적었고 노트
364가 ``팝업은 남의 학습이 빠질 때마다 잃는다''를 쟀다. 여기서 셋째 얼굴이
나온다 --- \textbf{팝업은 결정권까지 갖고 있다.}

\section{이번에는 안 버린다}

판 부호가 \textbf{9/12로 음수}다. 그것은 노트 351이 채택할 때 쓴 것
(팝업 11/12)과 같은 세기이고, 노트 361이 ``희석 반증''을 선언할 때 쓴 것
(6/12 대 8/12)보다 세다. \textbf{일하고 있는 축을 버리지 않는다.}

\textbf{그리고 여기서 폭을 고치지 않는다.} 고치면 이 결정이 규칙 덕에
바뀌는 모양이 되고, 그것은 결과를 보고 규칙을 고르는 것이다(노트 358이
같은 자리에서 반대로 겪었다 --- 그때는 내 사전 등록이 상시 규약과
어긋나서 상시 규약을 따랐다). \textbf{다음 노트에 미리 등록한다.}

\begin{quote}
\textbf{미리 등록(노트 367)} --- 미결정 폭을 고정 $0.010$ 에서
\emph{잰 값}으로 바꾼다: $|$판 차$| < 2 \times$ 짝SD \textbf{그리고}
부호가 $[4/12,\,8/12]$ 안일 때만 미결정으로 본다. 이 창의 여덟 실험에
소급해 적용해 어떤 결정이 바뀌는지 먼저 적고, \textbf{바뀌는 것이 있으면
그것을 밝힌 뒤} 규약을 바꾼다.
\end{quote}

\section{남은 것}

\begin{center}
\begin{tabular}{ll}
\hline
갈래 & 지금 아는 것 \\
\hline
\textbf{미결정 폭 고치기} & 위 사전 등록. 소급 적용부터 적는다 \\
만화 \texttt{media\_push} 다시 만들기 & AniList 외부 링크 · 예고편을 \\
 & 2,041건 다 긁는 중(950건 받음). 옛 값도 같이 \\
팝업 실측 41행 & 노트 360. 결정 다섯을 다 팝업이 내렸다 \\
세계애니 fill 도 자리 기준 & 노트 365. 같은 덫이 하나 더 \\
\hline
\end{tabular}
\end{center}

\end{document}
